\documentclass[11pt]{article}
\usepackage[margin=0.85in]{geometry}
\usepackage{lmodern}
\usepackage{microtype}
\usepackage{enumitem}
\usepackage{listings}
\usepackage[T1]{fontenc}

\lstset{
  basicstyle=\ttfamily\small,
  columns=fullflexible,
  keepspaces=true,
  showstringspaces=false,
  breaklines=true,
  frame=single,
  xleftmargin=0.4em,
  xrightmargin=0.4em,
  aboveskip=0.7em,
  belowskip=0.7em
}

\pagestyle{empty}
\setlength{\parskip}{0.45em}
\setlength{\parindent}{0pt}
\setlist[enumerate]{leftmargin=1.4em, itemsep=0.85em, topsep=0.4em}
\setlist[itemize]{leftmargin=1.3em, itemsep=0.25em, topsep=0.35em}

\begin{document}

\begin{center}
{\LARGE\bfseries Quiz 2 --- Sample 2}\\[0.35em]
DATA 412/612
\end{center}

\vspace{0.4em}

Write your answers in an R script. Use the tidyverse.
Load \texttt{nycflights13} and the \texttt{flights} data.
This is a second practice quiz. It is not a solution key.

\begin{enumerate}
\item In \texttt{flights}, for every \texttt{carrier}, find the number of
\texttt{NA}s in each column that is numeric and starts with \texttt{dep},
or that is character and contains \texttt{er} or \texttt{il}.

\item Restrict attention to winter flights (December, January, or February)
with \texttt{air\_time} greater than 180 minutes.
Find the carrier that operated the fewest of these flights.
Then, for every \texttt{origin}, count how many of that carrier's
winter long flights departed from that origin.

\item Among origin--destination pairs with at least 100 flights, which pair
has the highest average \texttt{distance}?
Report the origin, the destination, the number of flights, and the
average distance.

\item Which three months have the highest proportion of flights with
\texttt{arr\_delay} greater than 30 minutes?
Report the month, the number of flights, the number of those flights
that were delayed, and the proportion.
When you count delayed flights, ignore missing \texttt{arr\_delay}.
Keep every flight in the denominator.

\item Use \texttt{readr} parse functions to correct the type of each column
below. The column names are a hint.
\begin{itemize}
  \item \texttt{date} is year, month, day, with \texttt{..} between the pieces.
        The string \texttt{NA} means missing.
  \item \texttt{time} is hours, minutes, seconds, with \texttt{..} between the pieces.
  \item \texttt{num} is a number that may contain commas.
        The string \texttt{n/a} means missing.
  \item \texttt{log} is logical. The string \texttt{missing} means missing.
  \item \texttt{char} is stored as a factor. It should end up as character.
        The string \texttt{miss} means missing.
  \item \texttt{fac} should be a factor.
\end{itemize}

\begin{lstlisting}
df <- data.frame(
  date = c("2024..03..15", "2022..12..01", "NA"),
  time = c("09..15..00", "21..40..30", "07..05..05"),
  num  = c("1,200", "850", "n/a"),
  log  = c("TRUE", "missing", "FALSE"),
  char = as.factor(c("red", "blue", "miss")),
  fac  = c("X", "Y", "Z")
)
\end{lstlisting}
\end{enumerate}

\end{document}
